\documentclass[UTF8]{ctexart}
\usepackage[a4paper,margin=2.6cm]{geometry}
\begin{document}

\section*{第五部分：复杂度模型与 Universal Optimality}

为了讨论距离顺序问题的复杂度，论文引入了两个模型：comparison model 和 time model。

\subsection*{Comparison model}

在比较模型中，算法已经知道图的顶点和弧，但只能通过 oracle 访问弧长。算法允许做的关于弧长的操作，是比较两个弧长线性函数的大小。每做一次这样的比较，代价为 1。

在这个模型下，如果一个距离顺序算法对任意给定图所做的比较次数，都在该图所有可能弧长选择下所需最大比较次数的常数倍以内，那么称它在比较次数上是 universally optimal 的。

\subsection*{Time model}

在时间模型中，图通过顶点集、源点以及每个顶点的出弧邻接表给出。算法必须通过遍历邻接表来发现弧和弧长。访问邻接表中的第一条弧、访问某条弧之后的下一条弧，都各花费一个单位时间。如果列表为空或已经到末尾，就返回 null。

在这个模型里，算法仍然只能比较已经访问过的弧长所组成的线性函数。每次比较花费一个单位时间。作者指出，这个模型其实比真实机器更强，因为真实算法还要计算这些线性函数；因此，在这个模型下得到的下界会更有力。

如果一个距离顺序算法对任意给定图，在任意弧长和任意邻接表排列方式下，运行时间都在该图最坏表示所需时间的常数倍以内，那么称它在时间上是 universally optimal 的。

\subsection*{两个模型的区别}

比较模型比时间模型更强，但不如时间模型现实。一个极端例子是：图中有一条从源点出发的单路径，再加上一些从后面顶点指回前面顶点的弧。在比较模型里，距离顺序可能不需要任何比较就能确定；但在时间模型里，算法仍然可能需要花 \(\Omega(m)\) 时间检查输入中的弧。

\subsection*{这一节的作用}

第 5 节给后面的上下界定下评价标准。作者将在第 6 节证明：

\begin{itemize}
  \item 在时间模型中，距离顺序问题需要 \(\Omega(m+\log D)\) 时间；
  \item 在比较模型中，距离顺序问题需要 \(\Omega(F-n+1+\log D)\) 次比较。
\end{itemize}

其中，\(D\) 是图中可能的距离顺序数量，\(F\) 是任意距离顺序中最大可能的前向弧数量。

后面几节则会证明：合适实现的 Dijkstra 可以达到 \(O(m+\log D)\) 时间；经过增强后，还可以达到 \(O(F-n+1+\log D)\) 级别的比较次数。因此，Dijkstra 在这些模型下可以达到 universal optimality。

\end{document}
